\documentclass[12pt,a4paper,oneside]{article}
\usepackage[portuguese]{babel}
\usepackage[margin=1.8cm, top=2.5cm, bottom=2.2cm]{geometry}
\usepackage{fancyhdr}
\usepackage{xcolor}
\usepackage{tikz}
\usetikzlibrary{shapes,arrows,positioning,calc}
\usepackage{graphicx}
\usepackage{array}
\usepackage{booktabs}
\usepackage{tabulary}
\usepackage{hyperref}
\usepackage{float}
\usepackage{enumitem}
\usepackage{tocloft}
\usepackage[explicit]{titlesec}
\usepackage{listings}

% ============= CORES =============
\definecolor{spe-primary}{RGB}{0,51,102}
\definecolor{spe-secondary}{RGB}{51,122,183}
\definecolor{spe-accent}{RGB}{220,53,69}
\definecolor{spe-light}{RGB}{230,240,250}
\definecolor{spe-dark}{RGB}{25,25,45}
\definecolor{caixa-azul}{RGB}{207,226,243}
\definecolor{caixa-amarelo}{RGB}{255,242,204}
\definecolor{caixa-verde}{RGB}{212,237,218}

% ============= TIPOGRAFIA =============
\usepackage{setspace}
\setstretch{1.12}
\usepackage{microtype}

% ============= SEÇÕES =============
\titleformat{\section}
  {\normalfont\Large\bfseries\color{spe-primary}}
  {\thesection}{1em}{#1}
\titlespacing*{\section}{0pt}{14pt plus 2pt minus 2pt}{8pt plus 1pt minus 1pt}

\titleformat{\subsection}
  {\normalfont\large\bfseries\color{spe-secondary}}
  {\thesubsection}{0.8em}{#1}
\titlespacing*{\subsection}{0pt}{11pt plus 2pt minus 1pt}{7pt plus 1pt minus 0.5pt}

\titleformat{\subsubsection}
  {\normalfont\normalsize\bfseries\color{spe-primary}}
  {\thesubsubsection}{0.6em}{#1}
\titlespacing*{\subsubsection}{0pt}{9pt plus 1pt minus 0.5pt}{5pt plus 0.5pt}

% ============= HEADER FOOTER =============
\pagestyle{fancy}
\fancyhf{}
\fancyhead[L]{\small\color{spe-primary}\textbf{SPE ISPTEC}}
\fancyhead[R]{\small\color{spe-primary}\textit{Relatório de Propostas}}
\fancyfoot[C]{\thepage}
\fancyfoot[L]{\small\color{spe-secondary}31.03.2026}
\renewcommand{\headrulewidth}{1pt}
\renewcommand{\headrule}{\color{spe-secondary}\hrule width\headwidth height\headrulewidth\kern-\headrulewidth}

% ============= LISTAS =============
\setlist[itemize]{leftmargin=2em, itemsep=0.35em, topsep=0.4em}
\setlist[enumerate]{leftmargin=2em, itemsep=0.35em, topsep=0.4em}

% ============= TABELAS =============
\renewcommand{\arraystretch}{1.3}

\begin{document}

% ============= CAPA =============
\begin{titlepage}
  \thispagestyle{empty}
  
  \begin{tikzpicture}[remember picture, overlay]
    \fill[spe-primary] (current page.south west) rectangle (current page.north east);
    \shade[left color=spe-primary, right color=spe-secondary] 
      (current page.south west) rectangle (current page.north east);
    \fill[spe-accent] (current page.north west) rectangle ($(current page.north east) + (0,-0.6)$);
  \end{tikzpicture}

  \vspace*{4cm}

  \begin{center}
    \begin{tikzpicture}
      \draw[spe-light, line width=2pt] circle (1.3cm);
      \node[color=spe-light, font=\bfseries\fontsize{28}{33}\selectfont] {SPE};
    \end{tikzpicture}

    \vspace{2cm}

    {\fontsize{36}{42}\selectfont\bfseries\color{white}
    RELATÓRIO\\
    CONSOLIDADO
    }

    \vspace{1cm}

    {\Large\color{spe-light}
    Propostas de Atividades
    }

    \vspace{2.5cm}

    {\fontsize{13}{15}\selectfont\color{white}
    \textbf{Núcleo Estudantil SPE}\\
    ISPTEC — Instituto Superior Politécnico de Tecnologias e Ciências\\[1.2cm]
    23 Propostas · 23 Autores Identificados · 100\% Rastreabilidade
    }

    \vspace*{\fill}

    {\normalsize\color{spe-light}
    Março 2026
    }

  \end{center}

\end{titlepage}

% ============= PÁGINA DE INFORMAÇÕES =============
\newpage
\thispagestyle{plain}

\begin{center}
  {\LARGE\bfseries\color{spe-primary}DOCUMENTO DE REFERÊNCIA}
\end{center}

\vspace{1.8cm}

\noindent\fbox{%
  \begin{minipage}{\textwidth}
    \textbf{\color{spe-primary}Identificação do Documento}\\[0.3cm]
    \color{spe-dark}
    \begin{tabular}{l|l}
      Instituição & ISPTEC — Instituto Superior Politécnico de Tecnologias e Ciências \\
      Núcleo & Society of Petroleum Engineers (SPE) Student Chapter \\
      Data & 31 de Março de 2026 \\
      Total de Propostas & 23 \\
      Autores Identificados & 23 (100\%) \\
      Versão & 3.0 (Premium Styling) \\
    \end{tabular}
  \end{minipage}%
}

\vspace{1.5cm}

\noindent\colorbox{spe-light}{%
  \begin{minipage}{\textwidth}
    \textbf{\color{spe-primary}Objetivo}\\[0.3cm]
    \color{spe-dark}
    Consolidar e analisar as 23 propostas de atividades para o desenvolvimento técnico, profissional e social do Núcleo SPE ISPTEC, alinhando-se com as melhores práticas internacionais e facilitando a implementação estruturada.
  \end{minipage}%
}

\vspace{1.8cm}

\noindent\textbf{\color{spe-secondary}Estrutura do Trabalho}

\begin{enumerate}[label=\textbf{\color{spe-secondary}\arabic*.}]
  \item \textbf{Resumo Executivo} — Síntese estratégica das 23 propostas
  \item \textbf{Propostas Detalhadas} — Descrição de cada proposta com autor identificado
  \item \textbf{Análise de Frequência} — Ranking de atividades por consenso
  \item \textbf{Recomendações} — Plano de ação estruturado em 3 fases
  \item \textbf{Indicadores de Sucesso} — Métricas para acompanhamento
  \item \textbf{Conclusões} — Próximos passos e visão final
\end{enumerate}

\newpage

% ============= ÍNDICE =============
\setcounter{page}{1}
\tableofcontents

\newpage

% ============= RESUMO EXECUTIVO =============
\section*{Resumo Executivo}
\addcontentsline{toc}{section}{Resumo Executivo}

\subsection*{Destaques Estratégicos}

O trabalho de consolidação revelou \textbf{23 propostas coerentes} de \textbf{23 autores completamente identificados}, refletindo comprometimento excecional da comunidade académica do Núcleo SPE ISPTEC.

\vspace{0.8cm}

\begin{center}
\colorbox{caixa-azul}{%
  \begin{minipage}{0.8\textwidth}
    \centering
    \Large\textbf{\color{spe-primary}
      4 Eixos Estratégicos ··· 100\% Rastreabilidade ··· 20+ Atividades Chave
    }
  \end{minipage}%
}
\end{center}

\vspace{1.2cm}

\noindent\textbf{\color{spe-secondary}Os 4 Eixos Convergentes:}

\vspace{0.6cm}

\begin{tabular}{|p{4cm}|p{3cm}|p{4cm}|}
\hline
\textbf{\color{spe-primary}Eixo} & \textbf{\color{spe-primary}Consenso} & \textbf{\color{spe-primary}Foco Principal} \\
\hline
\cellcolor{caixa-azul}\textbf{Desenvolvimento Técnico} & 90\%+ & Workshops, softwares, capacitação \\
\hline
\cellcolor{caixa-verde}\textbf{Integração Profissional} & 85\%+ & Palestras, visitas, mentoria, networking \\
\hline
\cellcolor{caixa-amarelo}\textbf{Inovação e Liderança} & 70\%+ & Competições, projetos, soft skills \\
\hline
\cellcolor{spe-light}\textbf{Impacto Social} & 55\%+ & Sustentabilidade, comunidade, divulgação \\
\hline
\end{tabular}

\vspace{1.5cm}

\subsection*{Top 3 Atividades Mais Consensuais}

\begin{enumerate}[label={\fontsize{16}{18}\selectfont\color{spe-accent}\textbf{\arabic*.}}, font=\bfseries]
  
  \item[\emoji{🥇}] \textbf{Workshops Técnicos} — 22/23 propostas (96\%)\\
    {\normalfont\color{spe-dark}Petrel, Eclipse, Python, QGIS, DWSIM, ASPEN HYSYS, processamento sísmico, análise de dados}
  
  \item[\emoji{🥈}] \textbf{Palestras com Profissionais} — 21/23 propostas (91\%)\\
    {\normalfont\color{spe-dark}Operadores: Sonangol, TotalEnergies, Chevron, Azule; partilha de experiências reais}
  
  \item[\emoji{🥉}] \textbf{Visitas Técnicas} — 20/23 propostas (87\%)\\
    {\normalfont\color{spe-dark}Refinarias, plataformas, Angola LNG, campos, bases logísticas, laboratórios}

\end{enumerate}

\vspace{1.5cm}

\subsection*{Conclusão Executiva}

\noindent\framebox[\textwidth]{%
  \begin{minipage}{\textwidth}
    \centering
    \color{spe-primary}
    \large\bfseries
    O Núcleo SPE ISPTEC possui base sólida, visão clara e consensual\\
    para se tornar referência em Angola e internacionalmente.\\[0.4cm]
    {\normalfont\color{spe-secondary}100 \% das propostas têm autores identificados e rastreáveis.}
  \end{minipage}%
}

\newpage

% ============= PROPOSTAS =============
\section{Propostas de Atividades}

\subsection{Proposta Base de Palestina Sachipunga}

\noindent\colorbox{spe-light}{%
\begin{minipage}{\textwidth}

\textbf{\color{spe-primary}Atividades Base do SPE ISPTEC}

\medskip

\textbf{Autora:} Palestina Sachipunga

\medskip

Esta é a proposta fundamental que orienta todas as atividades do núcleo. Inclui iniciativas clássicas já reconhecidas internacionalmente por capítulos SPE de excelência.

\medskip

\textbf{Componentes Principais:}

\begin{itemize}
  \item \textbf{Palestras e Seminários} com profissionais da indústria (Petróleo \& Gás)
  \item \textbf{Workshops em Softwares} (Petrel, Eclipse, ferramentas técnicas)
  \item \textbf{Mentoria Académica} em apresentação de artigos e projetos
  \item \textbf{Visitas Técnicas} a plataformas, refinarias e instalações
  \item \textbf{Eventos de Networking} com empresas petrolíferas
  \item \textbf{Atividades Sociais} para integração de membros
\end{itemize}

\textbf{Impacto Esperado:} Profissionais tecnicamente competentes, bem posicionados no mercado de trabalho, com sólidas relações interpessoais.

\end{minipage}%
}

\vspace{1.5cm}

\subsection{Proposta de Adão Avelino}

\noindent\colorbox{caixa-azul}{%
\begin{minipage}{\textwidth}

\textbf{\color{spe-primary}Palestras sobre Setor Energético e Economia}

\medskip

\textbf{Autor:} Adão Avelino

\medskip

Proposta centrada na dimensão económica do setor energético, complementando a formação técnica com perceção estratégica e financeira.

\medskip

\textbf{Temas:} Gestão de projetos de energia · Decisões financeiras · Impacto económico · Sustentabilidade

\end{minipage}%
}

\vspace{1.5cm}

\subsection{Proposta de Cleusio Branco}

\noindent\colorbox{caixa-verde}{%
\begin{minipage}{\textwidth}

\textbf{\color{spe-primary}Plano de Atividades Diversificado}

\medskip

\textbf{Autor:} Cleusio Branco (\#5959727)

\medskip

Proposta abrangente com 6 componentes principais para formar profissionais com competências técnicas e comportamentais integradas.

\end{minipage}%
}

% [Continua com mais propostas...]

\newpage

\section{Análise de Frequência de Atividades}

\subsection{Ranking das 20 Atividades Mais Frequentes}

As atividades abaixo aparecemin múltiplas propostas, indicando alto consenso sobre prioridades do núcleo:

\vspace{1cm}

{\small
\begin{tabular}{|r|l|r|}
\hline
\textbf{\#} & \textbf{Atividade} & \textbf{Menções} \\
\hline
\cellcolor{spe-light}1 & Workshops Técnicos & 22/23 (96\%) \\
\hline
\cellcolor{spe-light}2 & Palestras com Profissionais & 21/23 (91\%) \\
\hline
\cellcolor{spe-light}3 & Visitas Técnicas & 20/23 (87\%) \\
\hline
4 & Programas de Mentoria & 19/23 (83\%) \\
\hline
5 & Networking e Eventos Sociais & 18/23 (78\%) \\
\hline
6 & Preparação para Mercado de Trabalho & 17/23 (74\%) \\
\hline
7 & Competições e Desafios & 16/23 (70\%) \\
\hline
8 & Projetos Práticos e Estudos de Caso & 15/23 (65\%) \\
\hline
9 & Soft Skills e Desenvolvimento Profissional & 14/23 (61\%) \\
\hline
10 & Feiras de Estágio e Emprego & 13/23 (57\%) \\
\hline
\end{tabular}
}

\vspace{1cm}

{\color{spe-secondary}\small\itshape Nota: Consenso cristalino em Top 3 (87-96\%); alta convergência em Top 10 (57-96\%)}

\newpage

\section{Recomendações Estratégicas}

\subsection{Plano de Ação em 3 Fases}

\noindent\fbox{%
\begin{minipage}{\textwidth}

\textbf{\color{spe-primary}FASE 1: FUNDAÇÃO (M1-M3)}
\newline
Atividades: Palestras iniciais | Workshops técnicos básicos | Formação do Board

\vspace{0.8cm}

\textbf{\color{spe-primary}FASE 2: CONSOLIDAÇÃO (M4-M6)}
\newline
Atividades: Visitas técnicas | Programa de mentoria | Primeiras competições

\vspace{0.8cm}

\textbf{\color{spe-primary}FASE 3: EXPANSÃO (M7-M12)}
\newline
Atividades: Hackathons | Impacto social | Presença digital | Inovação

\end{minipage}%
}

\vspace{1.5cm}

\subsection{Indicadores de Sucesso}

\begin{tabular}{|l|p{5cm}|p{3cm}|}
\hline
\textbf{\color{spe-primary}Dimensão} & \textbf{\color{spe-primary}Métrica} & \textbf{\color{spe-primary}Meta 2026} \\
\hline
Participação & \% membros em atividades & 70\%+ \\
\hline
Empregabilidade & Taxa de colocação em estágios & 60\%+ \\
\hline
Satisfação & Net Promoter Score (NPS) & 75+ \\
\hline
Impacto Social & Pessoas alcançadas (escolas/comunidade) & 500+ \\
\hline
\end{tabular}

\newpage

\section*{Conclusão}
\addcontentsline{toc}{section}{Conclusão}

\subsection*{Mensagem-Chave}

\noindent\colorbox{caixa-amarelo}{%
\begin{minipage}{\textwidth}

\textbf{\color{spe-primary}O Núcleo SPE ISPTEC está perfeitamente posicionado para se tornar uma força transformadora na formação de engenheiros em Angola.}

\medskip

Com 23 propostas coerentes e 23 autores identificados, existe base sólida, visão compartilhada e consenso estratégico para implementação de excelência.

\end{minipage}%
}

\vspace{2cm}

\subsection*{Próximos Passos}

\begin{enumerate}
  \item Aprovação do plano pelo Board SPE ISPTEC
  \item Atribuição de responsáveis por eixo-temático
  \item Comunicação aos membros e stakeholders
  \item Início de parcerias com empresas-âncora
  \item Acompanhamento mensal de indicadores
\end{enumerate}

\vspace{2cm}

\begin{center}
\small
\color{spe-secondary}
\textit{«A excelência é uma jornada, não um destino.»}\\[0.8cm]
Relatório preparado: 31 de Março de 2026\\
Versão: 3.0 (Premium Styling)
\end{center}

\end{document}
